\section{Compensation gives every Eisenstein sixfold determinant}

\begin{theorem}[HD-E6: an effective bound for every integral Eisenstein sixfold datum]
\label{HD-E6}
Let $(A,L)$ be an abelian sixfold of Weil type for
$K=\mathbf Q(\omega)$, with signature $(3,3)$ and integral
$\mathbf Z[\omega]$ action preserving the polarization. Let $\Delta$
be the product of the polarization elementary divisors. Once
D-G34-01--07 are applied with their checked proofs, every member of
this fixed datum has an effective integral three-cycle detecting its
Weil space, of degree
\[
 \delta_W(A,L)\le 2D_8(\Delta^2).
\]
The same formula is uniform on its whole Weil family, without a
restriction on the determinant class. Theorem HD-WQ then gives the
stated quadratic division bound for each normalized rational Weil class.
\end{theorem}
\begin{proof}
Use the Hermitian normalization of HECKE\_DIV SQ7,
$h(x,y)=E(x,\omega y)-\bar\omega E(x,y)$.
Its matrix on an integral Eisenstein basis has determinant $-\Delta$:
its determinant has sign $(-1)^3$, and its absolute value is the
product of the elementary divisors of $E$, by the determinant identity
in the cited proof.

Let $E_\omega$ have its degree-one principal class $\vartheta$, and put
\[
 T=E_\omega\times E_\omega,\quad
 a\cdot(x,y)=(ax,\bar a y),\quad
 \lambda=\vartheta_1+\Delta\vartheta_2.
\]
This is a polarized Weil surface with Hermitian matrix
$\operatorname{diag}(1,-\Delta)$ and polarization product $\Delta$.
The product $A\times T$, with class
$H=p_A^*c_1(L)+p_T^*\lambda$, has signature $(4,4)$, polarization
product $\Delta^2$, and Hermitian determinant $\Delta^2$.
That determinant is a norm from $K$, so the rational Hermitian space
is split, by Schoen (1998), Section~7, printed p.~332.
It is an integral maximal-order datum. D-G34-07 therefore supplies an
effective integral four-cycle $Z$ on $A\times T$ with
\[
 \deg_H Z\le D_8(\Delta^2),\qquad e_{W(A\times T)}[Z]\ne0.
\]
If needed, choose a reduced component with nonzero projection.

Here is a completely effective contraction. Denote by $F_1,F_2$
the two coordinate elliptic divisors on $T$, and by $\Gamma_1$ and
$\Gamma_\omega$ the graphs of $1$ and $\omega$. Their classes span
$\operatorname{NS}(T)_\mathbf Q$: the two diagonal Hermitian matrices
and the off-diagonal entries $1,\omega$ form a rational basis of
the Hermitian $2$-by-$2$ matrices over $K$. More specifically, the
classes
\[
 [F_1]+[F_2]-[\Gamma_1],\qquad
 [F_1]+[F_2]-[\Gamma_\omega]
\]
span $W_K(T)$, since they are precisely the two off-diagonal
Hermitian classes for the balanced action. The intersection pairing
on this plane is nondegenerate. This also follows directly from its
negative-definite Hermitian-matrix intersection form, or from the
surface Hodge index theorem.

The Kunneth summand of $[Z]$ in
$H^6(A)\otimes H^2(T)$ has a nonzero projection to
$W_K(A)\otimes W_K(T)$: the nonzero product Weil component is in
this summand, and over $\mathbf C$ its two lines are the same-embedding
determinant products. Contracting this tensor against all of
$W_K(T)$ cannot give zero. Thus at least one of the four effective
divisor classes just listed has
\[
 e_{W_K(A)}p_{A*}\bigl([Z]\,p_T^*[D]\bigr)\ne0.
\]
This conclusion needs no choice of a pure Weil cycle representing $Z$.
Indeed a nonzero contraction against a difference of the displayed
effective divisors forces a nonzero contraction against at least one
of them.

Translate that elliptic divisor generally on $T$. Such translates
meet every component of $Z$ properly: apply the elementary incidence
dimension theorem to $Z\to T$ and the transitive translation action
on $T$. Their intersections have effective integral multiplicities.
The pushforward to $A$ is therefore an effective integral three-cycle;
lower-dimensional images contribute zero. It detects by the preceding
cohomological calculation, since translation does not change its class.

Every one of these four divisors satisfies the nef inequality
$[D]\le2(\vartheta_1+\vartheta_2)\le2\lambda$.
For the coordinate divisors this is immediate. For either unit graph,
its Hermitian matrix has eigenvalues $0,2$, proving the first
inequality. Hence its pushed-forward degree is at most
\[
 \int_Z p_T^*[D]p_A^*c_1(L)^3
 \le 2\int_Z p_T^*\lambda\,p_A^*c_1(L)^3
 \le2\int_Z H^4
 \le2D_8(\Delta^2).
\]
All inequalities are intersections of nef classes with an effective
cycle. Passing to a detecting reduced component only decreases degree.
This proves the claimed integral bound. One could retain the positive
coefficient four in the binomial expansion of $H^4$ to sharpen the
last estimate, but no such sharpening is needed here.

The same fixed $T$ and polarization work for the entire fixed datum
of $A$, and the eightfold bound holds on every fibre, not just on the
open Hurwitz image. Thus the conclusion is uniform, including the
possibly nonsplit sixfold determinant. This is the effective version
of the compensation argument of Schoen (1998), Proposition~10,
printed pp.~332--333, with dimensions raised by two.
\end{proof}

\begin{corollary}[HD-E6b: a fixed nonmaximal order]
\label{HD-E6b}
If the sixfold datum above has an integral action only of
$R_f=\mathbf Z+f\mathbf Z[\omega]$, the bound
\[
 \delta_W(A,L)\le2D_8(f^{24}\Delta^2)
\]
holds uniformly for this fixed conductor and polarization datum.
\end{corollary}
\begin{proof}
Write $A=V/\Lambda$ and set $\Lambda'=f\mathbf Z[\omega]\Lambda$.
Then $f\Lambda\subset\Lambda'\subset\Lambda$, because $R_f$ acts
integrally. The new lattice is maximal-order stable and has index
$I=[\Lambda:\Lambda']\le f^{12}$. The identity on $V$ induces a
$K$-linear isogeny $p:A'=V/\Lambda'\to A$. Endow $A'$ with $p^*L$;
its polarization elementary-divisor product is $I\Delta$ by the
top intersection projection formula. Apply HD-E6 and push a detecting
effective three-cycle forward. Its degree against $L$ is exactly
its degree against $p^*L$, and its Weil projection stays nonzero.
The function $D_8(x)=(4G^8x)^8B_0$ is increasing for positive $x$.
Thus $2D_8((I\Delta)^2)\le2D_8(f^{24}\Delta^2)$, as required.
No claim independent of an unbounded conductor is made.
\end{proof}
